\begin{exercisechunk}
\item \textbf{Matrix inequalities in the interpolation proof.}
\label[exercise]{ex:l9-matrix-order}
\exerciseprofile{2}{2}{3}{\exproof\quad\excalculation}{%
Justify the spectral and matrix-order bounds used in the risk calculation.}
For symmetric matrices, $S\preceq R$ means
$x^\top Sx\le x^\top Rx$ for every vector $x$.
The matrix norm is the Euclidean operator norm
$\|M\|=\sup_{\|x\|=1}\|Mx\|$.
You may use the spectral theorem: every real symmetric matrix has an
orthonormal basis of eigenvectors.
\begin{enumerate}[(a),beginpenalty=10000]
\item Prove that $S\preceq R$ implies
$B^\top SB\preceq B^\top RB$ for any matrix $B$ of compatible size.
If $c_HI_n\preceq H\preceq C_HI_n$ with $c_H>0$, prove
\[
 C_H^{-1}I_n\preceq H^{-1}\preceq c_H^{-1}I_n,
 \qquad H^{-2}\preceq c_H^{-1}H^{-1}.
\]
Identify which inequalities follow by applying scalar inequalities to
the eigenvalues of the same symmetric matrix.

\item Suppose also $U^\top U\succeq(n/2)I_m$ for
$U\in\R^{n\times m}$. With
$A=U^\top H^{-1}U$ and $T=(I_m+A)^{-1}U^\top H^{-1}$,
prove
\[
 A\succeq\frac{n}{2C_H}I_m,\qquad
 \|(I_m+A)^{-1}\|\le\frac{2C_H}{n},
\]
and
\[
 TT^\top\preceq c_H^{-1}(I_m+A)^{-1}A(I_m+A)^{-1}
 \preceq\frac{2C_H}{c_Hn}I_m.
\]
Explain why it is legitimate to diagonalize $A$ and $(I_m+A)^{-1}$
simultaneously.

\item Set $K=H+UU^\top$, $\alpha=K^{-1}\mathbf Y$,
$H=D^{-1}VV^\top$, and $\widehat b=D^{-1}V^\top\alpha$.
Prove the full chain
\[
 \|\widehat b\|^2
 =D^{-1}\alpha^\top H\alpha
 \le D^{-1}\mathbf Y^\top K^{-1}\mathbf Y
 \le(c_HD)^{-1}\|\mathbf Y\|^2.
\]
\end{enumerate}
\begin{hints}
For (b), first write $TT^\top$ explicitly. For an eigenvalue $t>0$ of
$A$, use $t/(1+t)^2\le1/t$.
\end{hints}
\end{exercisechunk}
